\documentclass[11pt]{article}

\usepackage[margin=1in]{geometry}
\usepackage{amsmath,amssymb,amsthm}
\usepackage{mathtools}
\usepackage{enumitem}
\usepackage{hyperref}
\usepackage[dvipsnames]{xcolor}

\hypersetup{
  colorlinks=true,
  linkcolor=RoyalBlue,
  citecolor=RoyalBlue,
  urlcolor=RoyalBlue
}

\newcommand{\R}{\mathbb{R}}
\newcommand{\E}{\mathbb{E}}
\newcommand{\bvec}{\mathbf{b}}

\theoremstyle{plain}
\newtheorem{proposition}{Proposition}
\newtheorem{remark}{Remark}

\title{SRT Routing Hyperparameter Optimization\\[4pt]
\large Formulation, Per-Level Allocation, and Research Directions}
\author{}
\date{March 2026 --- Working Draft}

\begin{document}
\maketitle

% ------------------------------------------------------------------
\section{Problem Overview}
% ------------------------------------------------------------------

We study the problem of routing a natural-language query through a
Summary Representation Tree (SRT) to retrieve relevant code artifacts.
The router descends the tree level by level, selecting a subset of
children at each node for further expansion.  Three families of
hyperparameters govern the traversal:

\begin{itemize}[nosep]
  \item \textbf{Beam width} $b_l$: number of children retained at level~$l$.
  \item \textbf{Max depth} $d$: deepest level explored.
  \item \textbf{Token budget} $T$: total LLM tokens available for routing.
\end{itemize}

\noindent
The goal is to mathematically formalize the optimization over these
parameters, derive practical heuristics, and identify conditions under
which closed-form or near-closed-form solutions exist.

% ------------------------------------------------------------------
\section{Repository Model}
% ------------------------------------------------------------------

The repository is modeled as a rooted tree of depth~$H$ with
level-dependent branching factors $B_1, B_2, \ldots, B_H$.
Let
\[
  f_k \;=\; \Pr[\text{relevant item lies at depth } k],
  \qquad \sum_{k=1}^{H} f_k = 1.
\]
Empirically, $f_k$ is concentrated at moderate-to-large depths (most
queries target leaf-level source files, not top-level directories).
Estimating the support and shape of~$f_k$ on real repositories is a
key empirical prerequisite: if $f_k$ is concentrated on a narrow band
of depths, the optimal~$d$ is essentially determined, and the
remaining problem reduces to beam allocation under a budget.

Each relevant item at depth~$k$ carries a value $v_k \geq 0$.

% ------------------------------------------------------------------
\section{Routing Model}
% ------------------------------------------------------------------

At level~$l$, the router examines the children of each node in the
current beam and retains $b_l$ of them.  Define
\[
  r_l(b_l) \;=\; \Pr[\text{correct branch is retained at level } l
    \mid \text{beam width } b_l].
\]
The \emph{survival probability}---the chance that the correct path is
still in the beam after~$k$ levels---is
\[
  P_{\mathrm{survive}}(k;\, \bvec)
  \;=\; \prod_{l=1}^{k} r_l(b_l),
  \qquad \bvec = (b_1, \ldots, b_d).
\]

\paragraph{Parameterization.}
A natural saturating model for level-$l$ retention is
\begin{equation}\label{eq:retention}
  r_l(b_l) \;\approx\; 1 - \exp\!\bigl(-b_l / \alpha_l\bigr),
\end{equation}
where the \emph{difficulty parameter}
\[
  \alpha_l \;\propto\; B_l \cdot m_l
\]
scales with the branching factor $B_l$ and a per-level
ambiguity measure $m_l$ (e.g., average semantic similarity among
siblings).  This captures diminishing returns: doubling the beam from
$1 \to 2$ helps far more than $16 \to 32$.

% ------------------------------------------------------------------
\section{Expected Gain}
% ------------------------------------------------------------------

The expected gain from a routing policy $(\bvec, d)$ is
\begin{equation}\label{eq:gain}
  G(\bvec, d)
  \;=\; \sum_{k=1}^{d} f_k \, v_k \prod_{l=1}^{k} r_l(b_l).
\end{equation}

% ------------------------------------------------------------------
\section{Cost Model}
% ------------------------------------------------------------------

At level~$l$, the router processes $n_l(\bvec)$ nodes, each costing
$c_l$ tokens (the cost of reading that node's summary plus the LLM
prompt overhead).  Total cost:
\begin{equation}\label{eq:cost}
  C(\bvec, d) \;=\; \sum_{l=1}^{d} n_l(\bvec) \cdot c_l,
  \qquad
  C(\bvec, d) \;\leq\; T.
\end{equation}

For the LLM-oracle policy, $n_l$ is approximately $b_{l-1} \cdot B_l$
(every retained node at level $l{-}1$ expands $B_l$ children).
For the embedding policy, $c_l$ is negligible and the token budget is
effectively infinite---the binding constraint shifts to latency.

% ------------------------------------------------------------------
\section{Dilution Model}\label{sec:dilution}
% ------------------------------------------------------------------

Over-expansion does not merely waste tokens; it actively degrades
downstream ranking by forcing the LLM to discriminate among a larger
candidate set.  We model this via a \emph{dilution penalty}
$D(\bvec, d)$.

\paragraph{Proposed model.}
Let $\rho_l(\bvec)$ denote the fraction of nodes in the beam at
level~$l$ that are \emph{irrelevant} (i.e., not on any path to a
relevant leaf).  A natural penalty is
\begin{equation}\label{eq:dilution}
  D(\bvec, d)
  \;=\; \sum_{l=1}^{d}
    w_l \,\log\!\bigl(1 + n_l(\bvec)\bigr),
\end{equation}
where $w_l$ are level-dependent weights reflecting the cost of
ranking degradation at depth~$l$.  The logarithmic form captures the
empirical observation that ranking difficulty scales sublinearly in
candidate-set size.

An alternative formulation penalizes the irrelevant-to-relevant
ratio directly:
\[
  D'(\bvec, d)
  \;=\; \sum_{l=1}^{d} w_l \,\rho_l(\bvec).
\]
Both models are compatible with the objective below; the choice should
be guided by ablation experiments.

% ------------------------------------------------------------------
\section{Objective}
% ------------------------------------------------------------------

The combined objective is
\begin{equation}\label{eq:objective}
  U(\bvec, d)
  \;=\; G(\bvec, d)
  \;-\; \mu\, D(\bvec, d)
  \;-\; \lambda\, C(\bvec, d),
\end{equation}
and we seek
\[
  (\bvec^*, d^*)
  \;=\; \operatorname*{argmax}_{\bvec \in \mathbb{Z}_{>0}^d,\; d \leq H}
    U(\bvec, d)
  \qquad \text{subject to} \quad
  C(\bvec, d) \leq T.
\]

% ------------------------------------------------------------------
\section{Per-Level Beam Allocation}\label{sec:allocation}
% ------------------------------------------------------------------

Treating the beam widths as continuous and working in the log-survival
domain, the expected gain decomposes additively.  Taking logarithms of
the survival product in~\eqref{eq:gain} and applying the
parameterization~\eqref{eq:retention}:
\[
  \log P_{\mathrm{survive}}(k;\,\bvec)
  \;=\; \sum_{l=1}^{k}
    \log\!\bigl(1 - e^{-b_l/\alpha_l}\bigr).
\]

For fixed~$d$, the Lagrangian for maximizing $U$ subject to the budget
constraint is
\[
  \mathcal{L}(\bvec, \nu)
  \;=\; G(\bvec, d) - \mu\,D(\bvec,d)
        - \nu\bigl(C(\bvec,d) - T\bigr),
\]
where $\nu \geq 0$ is the multiplier on the token budget.  The KKT
conditions yield, for each level~$l$:
\begin{equation}\label{eq:kkt}
  \frac{\partial G}{\partial b_l}
  \;=\; \mu \,\frac{\partial D}{\partial b_l}
  \;+\; \nu \,\frac{\partial C}{\partial b_l}.
\end{equation}

\begin{proposition}[Water-filling heuristic]\label{prop:waterfill}
Under the exponential retention model~\eqref{eq:retention} and the
log-dilution model~\eqref{eq:dilution}, with $c_l$ approximately
constant across levels, the KKT conditions~\eqref{eq:kkt} prescribe
beam allocation proportional to the difficulty parameters:
\[
  b_l^* \;\propto\; \alpha_l
  \;=\; O(B_l \cdot m_l),
\]
with the constant of proportionality set by the token budget~$T$.
Harder levels (large branching factor, high ambiguity) receive wider
beams; easy levels are pruned aggressively.
\end{proposition}

\begin{proof}[Proof sketch]
Differentiating $G$ with respect to~$b_l$ and using the chain rule
through the product gives a marginal-gain term proportional to
$\frac{1}{\alpha_l} e^{-b_l/\alpha_l}$.  Setting the marginal gains
equal across levels (the water-filling condition) yields
$b_l / \alpha_l = \text{const}$, i.e., $b_l \propto \alpha_l$.
The budget constraint fixes the constant.
\end{proof}

\begin{remark}
This is the SRT analogue of water-filling in
information theory: the ``channel'' at each level has capacity governed
by $\alpha_l$, and the optimal policy allocates beam width (power) in
proportion to channel quality.
\end{remark}

% ------------------------------------------------------------------
\section{Adaptive Policies (Beyond Fixed Parameters)}
\label{sec:adaptive}
% ------------------------------------------------------------------

The fixed-parameter formulation optimizes $(\bvec, d)$ before the
query is issued.  In practice, the router observes LLM confidence
scores at each level, which provide real-time signal about whether the
beam is on the right track.

\paragraph{POMDP framing.}
At each level~$l$, the router observes a confidence vector
$\mathbf{s}_l \in [0,1]^{n_l}$ over the current beam.  An adaptive
policy $\pi$ maps the history
$(b_1, \mathbf{s}_1, \ldots, b_{l-1}, \mathbf{s}_{l-1})$ to the
next beam width $b_l$.  The optimal adaptive policy satisfies
\[
  \pi^* \;=\; \operatorname*{argmax}_\pi\;
    \E_{\mathbf{s}}\!\bigl[U(\bvec^{\pi}, d^{\pi})\bigr].
\]

The fixed-parameter optimum $(\bvec^*, d^*)$ provides a natural
baseline: any adaptive policy should beat the best fixed policy, and
the gap measures the value of adaptivity.

\paragraph{Simple adaptive heuristics.}
Before solving the full POMDP, two lightweight heuristics are worth
evaluating:

\begin{enumerate}[nosep]
  \item \textbf{Confidence-gated widening.}
    Start with the water-filling allocation $b_l^*$.  If the maximum
    confidence score at level~$l$ falls below a threshold~$\tau$,
    double the beam width (up to the budget).
  \item \textbf{Early termination.}
    If all confidence scores at level~$l$ exceed a high threshold
    $\tau_{\mathrm{high}}$, reduce $d$ to $l + 1$ (stop exploring
    deep) and reallocate the saved budget to wider beams at the
    current level.
\end{enumerate}

% ------------------------------------------------------------------
\section{Empirical Hypotheses}
% ------------------------------------------------------------------

\begin{enumerate}[nosep]
  \item \textbf{Optimal depth exists.}
    For any fixed budget~$T$, there is a depth $d^*$ beyond which
    additional levels reduce $U$ (the cost and dilution penalties
    outweigh the marginal gain from deeper items).
  \item \textbf{Beam width has diminishing returns.}
    $G$ is concave in each $b_l$; the first few beam slots provide
    most of the survival probability.
  \item \textbf{Over-exploration harms ranking.}
    Holding the token budget fixed, increasing beam width beyond
    $b_l^*$ degrades final retrieval precision (the dilution effect
    dominates).
  \item \textbf{Per-level allocation beats uniform allocation.}
    The water-filling policy $b_l \propto \alpha_l$ outperforms a
    uniform policy $b_l = b$ at the same total cost, especially when
    the $\alpha_l$ vary widely across levels.
  \item \textbf{Adaptive policies beat fixed policies.}
    Confidence-gated widening improves recall on ambiguous queries
    without degrading precision on easy queries.
\end{enumerate}

% ------------------------------------------------------------------
\section{Research Directions}
% ------------------------------------------------------------------

\paragraph{Estimation.}
Measure $f_k$, $\alpha_l$, and $c_l$ on real repositories.  The
depth distribution $f_k$ is the highest-priority empirical target:
if its support is narrow, the optimization simplifies dramatically.

\paragraph{Closed forms.}
The continuous relaxation of Proposition~\ref{prop:waterfill} should
yield explicit beam-allocation formulas via KKT conditions.  Extend
to the joint $(d, \bvec)$ problem by treating $d$ as a discrete outer
loop over a small range.

\paragraph{Dilution ablation.}
Compare the log-dilution model~\eqref{eq:dilution} against the
ratio model $D'$ and a pure cost-constraint baseline ($\mu = 0$) on
retrieval benchmarks.

\paragraph{Adaptive routing.}
Implement the confidence-gated and early-termination heuristics from
Section~\ref{sec:adaptive} and measure the gap between adaptive and
fixed-parameter performance.

\paragraph{Embedding vs.\ LLM oracle.}
The cost model differs qualitatively between the two routing policies
(token-limited vs.\ latency-limited).  Formulate a parallel
optimization for the embedding policy where the objective trades off
recall against traversal depth, with the token budget replaced by a
latency budget.

% ------------------------------------------------------------------
\section{Key Insights}
% ------------------------------------------------------------------

\begin{enumerate}[nosep]
  \item Hyperparameter tuning for SRT routing is a structured
    optimization over a tree with probabilistic routing and cost
    constraints.
  \item The optimal policy allocates beam width \emph{per level},
    not uniformly---yielding a water-filling solution where harder
    levels receive wider beams.
  \item The dilution penalty is not just a budget waste; it actively
    degrades ranking and must be modeled explicitly.
  \item Fixed-parameter optimization provides the baseline;
    adaptive policies that condition on observed confidence scores
    represent the natural next step.
\end{enumerate}

\end{document}
